\scanpage{37}
On est ramené à prouver que si \deleted{$S'$} on a un sous-schéma $S_0=V(J)$, $J^2=0$, et si $\widehat X_{\widehat S_0}$ est algé-brique, alors $\widehat X$ l'est aussi. On utilise la théorie du prolongement infinitésimal d'une structure lisse $X_0/S_0$ resp $\widehat X_0/\widehat S_0$, et :

$1^\circ)$ L'obstruction à prolonger $X_0$ (resp $\widehat X_0$) en un $X$ (resp. $\widehat X$) lisse sur $S$ (resp $\widehat S$) est dans ``le même'' groupe
\[
H^2(X_0,\unsure{\mathcal T|\mathfrak T}_{X_0/S_0}\otimes_{\mathcal O_{S_0}}J)\simeq H^{\unsure{2}}(\widehat X_0,\widehat{\unsure{\mathcal T|\mathfrak T}}_{\widehat X_0/\widehat S_0}\widehat\otimes_{\mathcal O_{\widehat S_0}}\widehat J
\]
et moyennant cette identification, l'obstruction est la $=$.

$2^\circ)$ L'indétermination pour remonter est dans ``le même'' groupe
\[
H^1(X_0,\unsure{\mathcal T|\mathfrak T}_{X_0/S_0}\otimes_{\mathcal O_{S_0}}J)\simeq H^{\unsure{1|2}}(\widehat X_0,\widehat{\unsure{\mathcal T|\mathfrak T}}_{\widehat X_0/\widehat S_0}\widehat\otimes_{\mathcal O_{\widehat S_0}}\widehat J).
\]

Il reste à prouver la

\uline{Proposition} Soit $S$ un \deleted{anneau local noeth complet} \inserted{schéma local \inserted{noeth} complet}, \deleted{$\dim>1$}, $(Z_i)_{i\in\mathbb N}$ une \deleted{\unsure{suite}} \inserted{famille dénombrable} de parties fermées \deleted{\unsure{strictes}} de $S$, telles que $Z=\bigcup_i Z_i$ contienne tt $x\in S$ avec $\dim\overline{\{x\}}\leq1$. Alors il existe une partie finie $J$ de $I$ telle que
\[
\bigcup_{i\in J}Z_i=S.
\]
\inserted{Soit $S=\operatorname{Spec}R$.}
Si $\dim S\leq1$, c'est trivial (alors \inserted{l'ens sous-jacent à} $S$ est fini, et l'hyp. signifie que $Z=S$). Si $\dim S\geq2$, on sait qu'il existe un anneau à val. discrète complet $V$, \deleted{\unsure{tel}} tel que $R$ soit fini sur $V[[T_1,\ldots,T_{\unsure{r}}]]$, \inserted{\unsure{réduit}} Notons qu'on est ramené \inserted{aussitôt} à prouver

\mspage{36}
